\documentclass[border=5mm]{standalone}
% ==============================================================================
% SETUP.TEX - CONFIGURACIÓN GLOBAL DEL PROYECTO
% ==============================================================================
% Este archivo centraliza todos los paquetes y estilos.
% Debe incluirse en el preámbulo del libro principal y de los archivos standalone.
% \PassOptionsToPackage{gray}{xcolor}
% --- 1. CODIFICACIÓN E IDIOMA ---
% fix-cm habilita las fuentes Computer Modern en tamaños arbitrarios (por debajo
% de 5 pt, entre otros). Debe cargarse antes que fontenc.
\usepackage{fix-cm}
\usepackage[utf8]{inputenc}
\usepackage[T1]{fontenc}
\usepackage[spanish, es-tabla]{babel}  
\usepackage{csquotes} % Recomendado para babel y citas

\usepackage{import}
% mode=image: Usa el PDF pre-compilado, nunca intenta recompilar.
% Debes compilar cada figura manualmente antes de compilar el libro.
\usepackage[mode=image, subpreambles=false]{standalone}

% --- 2. MATEMÁTICAS Y CIENCIAS ---
\usepackage{amsmath, amssymb, amsfonts, mathtools}
\usepackage{enumitem}

% LaTeX trae \DeclareMathSizes{5}{5}{5}{5}, así que a 5 pt (\tiny) los subíndices
% salen del mismo tamaño que la base: en las etiquetas de figuras el M de
% $\omega_M$ queda desproporcionado. Se restablece la proporción habitual.
\DeclareMathSizes{5}{5}{3.7}{3}

% --- 3. DISEÑO Y GEOMETRÍA --- 
% Unificamos los márgenes para todo el documento
\usepackage[left=2.5cm,right=2.5cm,top=3cm,bottom=3cm]{geometry}
\makeatletter
\ifstandalone % Evita que geometry dañe el recorte de las figuras bajándolas 3cm
  \@ifclasswith{standalone}{crop=false}{
    % Es un capítulo/sección: Dejamos que geometry actúe.
  }{
    % Es una figura: Neutralizamos geometry para que no las recorte mal.
    \geometry{pass}
  }
\fi
\makeatother
\usepackage{parskip} % Espaciado entre párrafos sin sangría
\usepackage{float}   % Para usar [H] en figura

% Ajustes de flotantes para evitar espacios verticales exagerados
\renewcommand{\topfraction}{0.95}      % Permite que las figuras ocupen hasta el 95% superior
\renewcommand{\bottomfraction}{0.95}   % Permite que las figuras ocupen hasta el 95% inferior
\renewcommand{\textfraction}{0.05}     % Permite hojas con tan solo un 5% de texto
\renewcommand{\floatpagefraction}{0.8} % Requiere que una página de solo figuras esté 80% llena
\setcounter{topnumber}{4}
\setcounter{bottomnumber}{4}
\setcounter{totalnumber}{10}

% --- 4. GRÁFICOS Y TABLAS ---
\usepackage{graphicx}
\usepackage{booktabs} % Tablas profesionales
\usepackage{caption}  % Control de captions y \captionof
\usepackage{tikz}
\usetikzlibrary{arrows.meta, calc, angles, quotes, babel, backgrounds}
\usepackage{pgfplots}
\usepgfplotslibrary{groupplots}
\usepgfplotslibrary{external}
\usepgfplotslibrary{patchplots}
\pgfplotsset{compat=1.18}
\usepackage[american, straightvoltages]{circuitikz}

% --- 5. COLORES Y CAJAS ---

\usepackage[table]{xcolor}
\usepackage{tcolorbox}

% Definición de colores corporativos/estilo
\definecolor{utnblue}{RGB}{0, 51, 102}
\definecolor{codegreen}{rgb}{0,0.6,0}
\definecolor{codegray}{rgb}{0.5,0.5,0.5}
\definecolor{codepurple}{rgb}{0.58,0,0.82}
\definecolor{backcolour}{rgb}{0.96,0.96,0.96}

% --- 6. ENTORNOS PERSONALIZADOS (CAJAS) ---

% Caja "Resultado" (Azul Corporativo) — Definiciones, fórmulas clave, resultados importantes.
% Uso: \begin{resultado}[Fórmula de Euler] ... \end{resultado}
\newtcolorbox{resultado}[1][]{
    colback=utnblue!5!white,
    colframe=utnblue,
    title=\textbf{#1},
    fonttitle=\bfseries,
    sharp corners,
    boxrule=0.5mm
}

% Caja "Nota" (Gris) — Advertencias, aplicaciones, contexto secundario.
% Uso: \begin{nota}[Cuidado con la calculadora] ... \end{nota}
% Uso: \begin{nota} ... \end{nota}  → título por defecto "Nota"
\newtcolorbox{nota}[1][Nota]{
    colback=codegray!5!white,
    colframe=codegray,
    title=\textbf{#1},
    fonttitle=\bfseries,
    sharp corners,
    boxrule=0.5mm
}

% Caja inline para resaltar un valor y enlazarlo con su otra aparición en una tabla
% o en una cuenta: mismo color, mismo valor.
% Uso: \marcavalor{$4$}  o  \marcavalor[codegreen]{$4$}
\newtcbox{\marcavalor}[1][utnblue]{
    on line,
    boxsep=1pt, left=2pt, right=2pt, top=1pt, bottom=1pt,
    colback=#1!10!white,
    colframe=#1,
    boxrule=0.4pt,
    arc=2pt
}

% --- 7. CÓDIGO FUENTE (LISTINGS) ---
\usepackage{listings}

\lstdefinestyle{miestilo}{
    backgroundcolor=\color{backcolour},   
    commentstyle=\color{codegreen},
    keywordstyle=\color{magenta},
    numberstyle=\tiny\color{codegray},
    stringstyle=\color{codepurple},
    basicstyle=\ttfamily\footnotesize,
    breakatwhitespace=false,         
    breaklines=true,                 
    captionpos=b,                    
    keepspaces=true,                 
    numbers=left,                    
    numbersep=5pt,                  
    showspaces=false,                
    showstringspaces=false,
    showtabs=false,                  
    tabsize=2,
    frame=single,                    
    rulecolor=\color{codegray}
}

\lstset{style=miestilo}
% Soporte para tildes y ñ en bloques de código
\lstset{literate={ñ}{{\~n}}1 {á}{{\'a}}1 {é}{{\'e}}1 {í}{{\'i}}1 {ó}{{\'o}}1 {ú}{{\'u}}1}

% Operadores matemáticos personalizados

% Vector en sentido discreto (lista finita de coordenadas): negrita + flecha.
% Uso: \vect{v}, \vect{e}_k, \vect{0}.
\newcommand{\vect}[1]{\vec{\mathbf{#1}}}

% Fecha de versión y comando para capítulos standalone
\providecommand{\verdate}{\the\year.\ifnum\month<10 0\fi\the\month.\ifnum\day<10 0\fi\the\day}
\newcommand{\chapterversion}{%
    \vspace{-1.5cm}\begin{flushright}\small\textit{\color{codegray}Versión~\verdate}\end{flushright}\vspace{0.5cm}%
}

% --- 8. HIPERVÍNCULOS (DEBE SER EL ÚLTIMO PAQUETE) ---
\usepackage[colorlinks=true, linkcolor=utnblue, urlcolor=utnblue, citecolor=utnblue]{hyperref}

% --- 9. REFERENCIAS CRUZADAS ENTRE CAPÍTULOS STANDALONE ---
% Cuando se compila un capítulo standalone, xr-hyper lee los .aux de los
% otros capítulos (si existen) para resolver \ref{} a labels externos.
% En la compilación del libro completo este bloque no se activa.
\ifstandalone
  \usepackage{xr-hyper}
  \makeatletter
  \newcommand{\xrchap}[1]{%
    \edef\xrc@self{\detokenize{Capitulo#1}}%
    \edef\xrc@job{\jobname}%
    \ifx\xrc@self\xrc@job\else
      \IfFileExists{../#1capitulo/Capitulo#1.aux}%
        {\externaldocument{../#1capitulo/Capitulo#1}}{}%
    \fi
  }
  \makeatother
  \xrchap{01}\xrchap{02}\xrchap{03}\xrchap{04}
  \xrchap{05}\xrchap{06}\xrchap{07}\xrchap{08}
  \xrchap{09}\xrchap{10}\xrchap{11}\xrchap{12}
  \xrchap{13}
\fi
% \PassOptionsToPackage{gray}{xcolor}

% fig_antialiasing
% Espectro realista de una señal del mundo: un lóbulo principal que decae con la
% frecuencia (forma de pasabajos de primer orden) más una cola tenue de nivel bajo
% que se extiende indefinidamente (ruido, contenido de alta frecuencia). NO está
% limitado en banda: la cola nunca llega a cero. Ese es el punto de la figura.
%   X(w) = 1/(1+(w/0.9)^2) + 0.05/(1+(w/5)^2)
% Muestreo: ws = 6 (=> ws/2 = 3). Corte del filtro: wM = 2, de modo que
% ws = 6 > 2 wM = 4, con separación visible entre copias (lo que se pide en la
% prosa: corte por debajo de Nyquist / muestreo algo más rápido que el mínimo).
% (a) el espectro de la señal, antes de cualquier operación: hay contenido más allá
%     de ws/2.
% (b) sin filtro: copias en 0, +-6 (rojo punteado); la suma (azul) queda por
%     encima de la copia central dentro de la banda base. El exceso (rojo pleno) es
%     el alias: las colas de las vecinas sumadas al contenido útil.
% (c) con filtro: el pasabajos recorta en wM=2. Lo descartado ---la cola gris--- es
%     de nivel bajo frente al lóbulo que se conserva. Copias limpias y separadas.

\begin{document}
\begin{tikzpicture}[>={Latex[length=4pt]}, font=\scriptsize, x=0.50cm, y=1.05cm,
    declare function={
        X(\x)  = 1/(1+(\x/0.9)^2) + 0.05/(1+(\x/5)^2);
        S(\x)  = X(\x) + X(\x-6) + X(\x+6);
    },
    tri/.style={utnblue, thick, fill=utnblue!12},
    tricopy/.style={utnblue!70, thick, fill=utnblue!6},
    suma/.style={utnblue, very thick},
    copia/.style={red!75!black, semithick, dash pattern=on 2pt off 1.6pt},
    desc/.style={gray!60, thin, fill=gray!22},
    lp/.style={red!80!black, thick, dotted},
    nyq/.style={gray!70, thin, dashed},
    ax/.style={->, thin, black},
]
    % ===================== (a) el espectro de la señal =====================
    \begin{scope}[shift={(0,0)}]
        \node at (-9.6,0.62) {(a)};
        \begin{scope}
            \clip (-8.5,-0.06) rectangle (8.5,1.42);
            \draw[tri] (-8.5,0)
                -- plot[domain=-8.5:8.5, samples=220, variable=\x] (\x,{X(\x)})
                -- (8.5,0) -- cycle;
        \end{scope}
        \draw[ax] (-8.9,0) -- (9.1,0) node[right] {$\omega$};
        \draw[ax] (0,-0.06) -- (0,1.42);
        \node[below, font=\tiny] at (0,-0.06) {$0$};
        \node[utnblue, anchor=west, font=\tiny] at (0.9,1.06) {$X(j\omega)$};
        % banda base del muestreo que se va a usar
        \draw[nyq] (3,-0.06) -- (3,1.30);
        \draw[nyq] (-3,-0.06) -- (-3,1.30);
        \node[gray!70, anchor=south, font=\tiny] at (3,1.30) {$\tfrac{\omega_s}{2}$};
        % la cola no se anula
        \node[gray!60!black, anchor=south, font=\tiny] at (5.6,0.36) {cola};
        \draw[gray!60!black, thin, -{Latex[length=3pt]}] (5.6,0.34) -- (5.6,0.09);
    \end{scope}

    % ===================== (b) sin filtro =====================
    \begin{scope}[shift={(0,-2.45)}]
        \node at (-9.6,0.62) {(b)};
        \begin{scope}
            \clip (-8.5,-0.06) rectangle (8.5,1.42);
            % suma de todas las copias
            \fill[utnblue!12] (-8.5,0)
                -- plot[domain=-8.5:8.5, samples=220, variable=\x] (\x,{S(\x)})
                -- (8.5,0) -- cycle;
            % exceso dentro de la banda base = alias
            \fill[red!25]
                   plot[domain=-3:3, samples=100, variable=\x] (\x,{S(\x)})
                -- plot[domain=3:-3, samples=100, variable=\x] (\x,{X(\x)}) -- cycle;
            % cada copia por separado
            \foreach \c in {-6,0,6} {
                \draw[copia] plot[domain=-8.5:8.5, samples=220, variable=\x]
                    (\x,{X(\x-\c)});
            }
            % contorno de la suma
            \draw[suma] plot[domain=-8.5:8.5, samples=220, variable=\x] (\x,{S(\x)});
        \end{scope}
        \draw[ax] (-8.9,0) -- (9.1,0) node[right] {$\omega$};
        \draw[ax] (0,-0.06) -- (0,1.42);
        \node[below, font=\tiny] at (0,-0.06) {$0$};
        % banda base
        \draw[nyq] (3,-0.06) -- (3,1.30);
        \draw[nyq] (-3,-0.06) -- (-3,1.30);
        \node[gray!70, anchor=south, font=\tiny] at (3,1.30) {$\tfrac{\omega_s}{2}$};
        \draw[thin] (6,0.05) -- (6,-0.05) node[below, font=\tiny] {$\omega_s$};
        \draw[thin] (-6,0.05) -- (-6,-0.05) node[below, font=\tiny] {$-\omega_s$};
        % etiquetas de los trazos
        \node[utnblue, anchor=south, font=\tiny] at (7.3,0.78) {suma};
        \draw[utnblue, thin, -{Latex[length=3pt]}] (7.3,0.76) -- (7.3,0.20);
        \node[red!75!black, anchor=south, font=\tiny] at (-7.3,0.78) {copias};
        \draw[red!75!black, thin, -{Latex[length=3pt]}] (-7.3,0.76) -- (-7.3,0.14);
        \node[red!75!black, anchor=south, font=\tiny] at (1.5,0.68) {alias};
        \draw[red!75!black, thin, -{Latex[length=3pt]}] (1.7,0.68) -- (2.2,0.27);
    \end{scope}

    % ===================== (c) con filtro =====================
    \begin{scope}[shift={(0,-4.90)}]
        \node at (-9.6,0.62) {(c)};
        \begin{scope}
            \clip (-8.5,-0.06) rectangle (8.5,1.42);
            % lo que el filtro descarta
            \fill[desc] (2,0) -- plot[domain=2:4, samples=80, variable=\x]
                    (\x,{X(\x)}) -- (4,0) -- cycle;
            \fill[desc] (-2,0) -- plot[domain=-2:-4, samples=80, variable=\x]
                    (\x,{X(\x)}) -- (-4,0) -- cycle;
            % copias vecinas del espectro filtrado
            \foreach \c in {-6,6} {
                \pgfmathsetmacro{\la}{\c-2}
                \pgfmathsetmacro{\lb}{\c+2}
                \draw[tricopy] (\la,0)
                    -- plot[domain=\la:\lb, samples=100, variable=\x] (\x,{X(\x-\c)})
                    -- (\lb,0);
            }
            % copia central
            \draw[tri] (-2,0)
                -- plot[domain=-2:2, samples=100, variable=\x] (\x,{X(\x)})
                -- (2,0);
        \end{scope}
        \draw[ax] (-8.9,0) -- (9.1,0) node[right] {$\omega$};
        \draw[ax] (0,-0.06) -- (0,1.42);
        \node[below, font=\tiny] at (0,-0.06) {$0$};
        % respuesta del filtro antialiasing
        \draw[lp] (-2,0) -- (-2,1.25) -- (2,1.25) -- (2,0);
        \node[red!80!black, anchor=south, font=\tiny] at (0,1.26) {filtro};
        % banda base y marcas
        \draw[nyq] (3,-0.06) -- (3,1.30);
        \draw[nyq] (-3,-0.06) -- (-3,1.30);
        \node[gray!70, anchor=south, font=\tiny] at (3,1.30) {$\tfrac{\omega_s}{2}$};
        \draw[thin] (2,0.05) -- (2,-0.05) node[below, font=\tiny] {$\omega_M$};
        \draw[thin] (6,0.05) -- (6,-0.05) node[below, font=\tiny] {$\omega_s$};
        \draw[thin] (-6,0.05) -- (-6,-0.05) node[below, font=\tiny] {$-\omega_s$};
        % separación entre copias
        \draw[<->, gray!65] (2,0.55) -- (4,0.55);
        \node[gray!65, anchor=south, font=\tiny, fill=white, inner sep=1pt]
            at (3,0.58) {separación};
        % lo descartado
        \node[gray!60!black, anchor=south, font=\tiny, fill=white, inner sep=1pt]
            at (-5.2,0.42) {descartado};
        \draw[gray!60!black, thin, -{Latex[length=3pt]}] (-4.6,0.42) -- (-3.4,0.10);
    \end{scope}
\end{tikzpicture}
\end{document}
